% \documentclass[twocolumn]{article}

\documentclass{article}

\usepackage[margin=1in]{geometry}
\usepackage[en-US]{datetime2} 
\usepackage[utf8]{inputenc}
\usepackage{graphicx}
\usepackage{amsmath}
\usepackage{hyperref}
\usepackage{tikz}
\usetikzlibrary{positioning}

\raggedright

\title{Domain Adaptation and Reasoning Frameworks in Language Models: A Controlled Experiment with Historical Cosmology}
\author{Francesco De Bernardis}
% \date{}

\begin{document}
\DTMlangsetup{showdayofmonth=false}


% \twocolumn[
% \begin{@twocolumnfalse}

\maketitle

\begin{abstract}
We investigate how domain adaptation reshapes explanatory behavior in language models using historical cosmology as a controlled setting. In Phase~1, we train a small language model from scratch on a pre-Copernican corpus from which explicit heliocentric references were removed, and evaluate whether Earth-motion or heliocentric continuations nevertheless emerge. In Phase~2, we fine-tune a larger pretrained model using QLoRA on the same corpus in order to study how adaptation modifies explanatory framing and cosmological stance.

Model outputs are evaluated using an LLM-as-judge framework that labels both cosmological stance (geocentric, heliocentric, or ambiguous) and explanatory frame (premodern versus modern). In the constrained setting of Phase~1, the smaller models occasionally generate local Earth-motion continuations, but these remain globally unstable and insufficient to support coherent cosmological reasoning. In Phase~2, fine-tuning induces a large and statistically significant shift toward premodern explanatory framing, while the conditional cosmological stance distributions remain comparatively stable within those frames. As a result, increases in geocentric outputs arise primarily from redistribution over explanatory regimes rather than from direct modification of stance.

These results suggest that domain adaptation may primarily reshape the explanatory regimes from which continuations are generated, with changes in stance emerging secondarily from those shifts.

\end{abstract}
\vspace{1em}

% \end{@twocolumnfalse}
]



\section{Introduction}

Can a language model generate heliocentric ideas despite being trained exclusively on geocentric texts?

We investigate this question in two regimes: small models trained from scratch under strict data constraints, and large pretrained models adapted via parameter-efficient fine-tuning (QLoRA). These settings allow us to probe both the emergence of new ideas and the selection between competing explanatory frameworks.
% \input{sections/data}
\section{Experimental Design}\label{sec:exp_design}


\subsection{Phase 1: Small-model training}

In Phase 1, we train a 110M-parameter GPT model on a filtered general-language corpus with astronomical content removed, followed by fine-tuning on a smaller pre-Copernican astronomy corpus.

The purpose of the general corpus is to provide broad exposure to English syntax, vocabulary, and discourse structure while minimizing direct exposure to modern astronomical concepts. The corpus was constructed from the Project Gutenberg archive \cite{gutenberg}, selecting a total of 2,851 documents using metadata-based filtering to exclude texts explicitly related to astronomy or science.

Because constructing a sufficiently large training corpus exclusively from Medieval or premodern texts was not feasible using publicly available sources, the general corpus also includes relatively recent literary and historical works. However, metadata filtering alone was insufficient to guarantee the removal of modern astronomical knowledge. For example, otherwise unrelated texts could still contain references to heliocentrism, planetary motion, or modern cosmology.

To reduce this contamination, the selected documents were further processed through keyword- and pattern-based filtering designed to remove passages referring to concepts such as Earth's orbit, heliocentrism, Copernicus, Galileo, and related modern astronomical ideas, while preserving general linguistic structure and non-astronomical content from literature, philosophy, and history. While complete removal of all indirect astronomical references cannot be guaranteed, the filtering procedure substantially reduced explicit exposure to heliocentric concepts and modern astronomical explanations.

The astronomy corpus is substantially smaller than the general corpus due to the limited availability of publicly accessible English translations of pre-Copernican texts. The corpus combines classical, late antique, and Medieval works containing geocentric astronomical reasoning, cosmological discussion, and premodern natural philosophy. Representative texts include Sacrobosco’s \textit{Sphaera Mundi}, Ptolemy’s \textit{Almagest}, Plato’s \textit{Timaeus}, Aristotle’s \textit{De Caelo}, Cleomedes’s \textit{On the Heavens}, and Peuerbach’s \textit{Theoricae Novae Planetarum}. Some works are explicitly astronomical, while others embed geocentric cosmology within broader philosophical or theological discussions. For a more comprehensive list see Appendix~\ref{appendix:astronomy_corpus}. Because many modern translations include editorial notes, annotations, or footnotes containing references to contemporary astronomy, the same pattern-based filtering procedure was applied to the astronomy corpus in order to remove translator commentary and other potential sources of modern astronomical leakage.






\subsection{Phase 2: QLoRA adaptation}

We adapt Qwen2.5-7B using QLoRA on the same astronomy corpus. Evaluation is performed using paired prompts and LLM-as-judge labeling.

Although the premodern explanatory frame is defined using historical astronomical machinery associated with geocentric models, it does not imply a geocentric stance. In practice, most premodern-frame outputs remain ambiguous or non-committal, making the mapping between frame and stance an empirical property rather than a definitional constraint.


\subsection{LLM-as-judge}
We generate 1680 prompts per model: 28 prompts, 4 categories, 15 samples

% define labels
%earth-motion, expliticit earth-motion, proto-heliocentirc, geocentric, ambiguous
%Geocentric classifications do not necessarily correspond to explicit assertions that the Earth is at the center of the universe. In many cases they reflect outputs that are consistent with a geocentric explanatory frame, typically characterized by the absence of Earth-motion and the use of premodern astronomical assumptions, and the recourse to pre-modern concepts, such has epycicles. 
% In this work, “geocentric” outputs should be understood as outputs consistent with a geocentric explanatory framework, rather than explicit doctrinal assertions about the Earth's centrality.
%refined_stance: stricter than stance
\section{Results}

\subsection{Phase 1}
Both the general English pre-trained model A and its astronomy fine-tuned version (model B), generate heliocentric-like content at low rates $(4\% - 8\%)$, suggesting some pattern recombination occurs even in a small 110M parameter models trained on geocentric texts. As any heliocentric content was completely removed from the training corpus, the emergence of this content occurs solely through pattern recombination. However, the quality of the generated output remains mediocre, with the overwhelming majority ($\sim 94\%$) of the output being labeled with the intermediate score of $1$, representing 
only "partly coherent" output according to our judge manifest. This is not surprising and a known limit of small models (see for example \cite{eldan2023tinystoriessmalllanguagemodels}).

The more interesting result is that Model A, despite being trained on an astronomy-free corpus, shows higher heliocentric content than Model B. Specifically, higher Earth-motion mentions rates ($8.3\%$ for model A vs $5.7\%$ for model B) and higher explicit Earth-motion mentions ($4.2\%$ (A) vs $3.3\%$ (B))
and proto-heliocentric content ($5.1\%$ (A) vs $4.0\%$ (B)). At the same time, surprisingly, model B also shows a lower rate of geocentric content ($3.3\%$ for A vs $1.9\%$ for B). This apparent contradiction is explained by the observation that the ambiguity of Model B statements increased compared to Model A, from $41\%$ (A) to $54\%$. In other words, Model B did not become more heliocentric nor more geocentric after the astronomy fine-tuning. Instead, the fine-tuning increased ambiguity rather than strengthening geocentric doctrine. 

In order to assess the statistical significance of these results, we conducted a prompt-level paired permutation test \cite{good2005permutation}. Since each prompt was sampled $15$ times per model, the 1680 generations per model are not independent observations. Some prompts are intrinsically more likely than others to elicit Earth-motion or ambiguity, and those within-prompt samples are correlated. For each of the $112$ prompts, the average rate $r$ for a given label was computed separately for A and B. This yields one paired difference per prompt, for each label $\ell$:

$$d^{\ell}_{AB,prompt} = r^{\ell}_{A,prompt} - r^{\ell}_{B,prompt}$$

The permutation test goal is to compare the mean of those paired differences with the expected mean under the null hypothesis that A and B are exchangeable at the prompt level. To accomplish this, the test flipped the sign of each paired difference $10000$ times, at random, and recalculated the mean. The two-sided $p$-value is the proportion of random sign-flips whose absolute mean difference is at least as large as the observed one. The advantage of this method is that it is a non-parametric test, requiring no assumption about the distribution of the differences. Table \ref{tab:phase1_results} shows the most relevant values. The drop in both Earth-motion mentions and Geocentric mentions is significant, but the most statistically significant results is the increase in Ambiguous statements. For a detailed analysis by prompt category see Appendix \ref{Appendix A}.
%Continue with % and modestly reduces low-bar Earth-motion mention and geocentric commitment.
% This distinction matters. The fine-tune seems to be reshaping the discourse mode more strongly than it is decisively suppressing Earth-motion assertions.

\begin{table}[t]
\centering
\caption{Comparison of heliocentric and ambiguity-related behaviors between Model A (general pre-training only) and Model B (astronomy fine-tuned). Values are averaged over paired prompts.}
\label{tab:phase1_results}
\begin{tabular}{lccc}
\hline
Metric & A & B & $p$-value \\
\hline
Earth-motion mention & 8.3\% & 5.7\% & 0.0226$^{**}$ \\
Explicit Earth-motion & 4.2\% & 3.3\% & 0.0954 \\
Proto-heliocentric & 5.1\% & 4.0\% & 0.1500 \\
Geocentric & 3.3\% & 1.9\% & 0.0136$^{**}$ \\
Ambiguous & 41.0\% & 54.8\% & $<10^{-4\,***}$ \\
\hline
\end{tabular}
\vspace{0.5em}
\footnotesize{
$^{**} p < 0.05$, \quad
$^{***} p < 0.001$
}
\end{table}






% \begin{table}[t]
% \centering
% \caption{Frequency of selected hedging expressions in generated outputs. Model B (astronomy fine-tuned) exhibits substantially more scholastic hedging than Model A.}
% \label{tab:hedging}
% \begin{tabular}{lcc}
% \hline
% Phrase & Model A & Model B \\
% \hline
% it seems & 1.7\% & 6.0\% \\
% according to & 6.7\% & 11.0\% \\
% it would seem & 1.7\% & 3.2\% \\
% any hedge phrase & 11.6\% & 20.2\% \\
% \hline
% \end{tabular}
% \end{table}




\subsection{Phase 2}

QLoRA induces a strong shift toward premodern explanatory framing, with approximately 25\% of paired samples flipping from modern to premodern regimes.
\section{Discussion}

The results suggest that fine-tuning primarily alters the selection of explanatory frameworks rather than directly encoding beliefs. This raises the possibility that different reasoning regimes correspond to distinct regions in the model’s activation space.
%Continue with % and modestly reduces low-bar Earth-motion mention and geocentric commitment.
% This distinction matters. The fine-tune seems to be reshaping the discourse mode more strongly than it is decisively suppressing Earth-motion assertions.
\section*{Acknowledgements}
The author is grateful to James Triveri for encouragement and for maintaining continued interest in this project.
\appendix
\section{Appendix A}\label{Appendix A}
\begin{table*}[t]
\centering
\caption{Category-level comparison between Model A and Model B across evaluation metrics.}
\label{tab:phase1_category_results}
\begin{tabular}{llcccc}
\hline
Category & Metric & A & B & Difference (B-A) & $p$-value \\
\hline
Astro & Earth-motion mention & 3.8\% & 1.9\% & -1.9\% & 0.3089 \\
Astro & Proto-heliocentric & 2.9\% & 1.0\% & -1.9\% & 0.1021 \\
Astro & Ambiguous & 60.0\% & 84.3\% & +24.3\% & $<10^{-4}$ \\
\hline
Declarative & Earth-motion mention & 2.1\% & 1.9\% & -0.2\% & 0.9959 \\
Declarative & Proto-heliocentric & 1.2\% & 1.2\% & 0.0\% & 1.0000 \\
Declarative & Ambiguous & 41.7\% & 59.1\% & +17.4\% & 0.0002 \\
\hline
General & Earth-motion mention & 0.0\% & 0.2\% & +0.2\% & 1.0000 \\
General & Ambiguous & 1.4\% & 5.2\% & +3.8\% & 0.0034 \\
\hline
Questions & Earth-motion mention & 27.4\% & 18.8\% & -8.6\% & 0.0451 \\
Questions & Proto-heliocentric & 16.2\% & 13.6\% & -2.6\% & 0.3172 \\
Questions & Ambiguous & 60.7\% & 70.5\% & +9.8\% & 0.0118 \\
\hline
\end{tabular}
\end{table*}

\begin{table*}[t]
\centering
\caption{Frequency of scholastic hedging expressions across generated outputs.}
\label{tab:full_hedging}
\begin{tabular}{l l r r r}
\hline
Model & Phrase & Total Texts & Phrase Count & Rate \\
\hline
A & some say & 1680 & 4 & 0.24\% \\
A & it may be understood & 1680 & 0 & 0.00\% \\
A & it seems & 1680 & 28 & 1.67\% \\
A & according to & 1680 & 113 & 6.73\% \\
A & if one grants & 1680 & 15 & 0.89\% \\
A & it may be supposed & 1680 & 0 & 0.00\% \\
A & some have held & 1680 & 0 & 0.00\% \\
A & it would seem & 1680 & 28 & 1.67\% \\
A & it may be asked & 1680 & 0 & 0.00\% \\
A & one may suppose & 1680 & 15 & 0.89\% \\
A & any hedge phrase & 1680 & 195 & 11.61\% \\
\hline
B & some say & 1680 & 4 & 0.24\% \\
B & it may be understood & 1680 & 0 & 0.00\% \\
B & it seems & 1680 & 100 & 5.95\% \\
B & according to & 1680 & 185 & 11.01\% \\
B & if one grants & 1680 & 15 & 0.89\% \\
B & it may be supposed & 1680 & 0 & 0.00\% \\
B & some have held & 1680 & 0 & 0.00\% \\
B & it would seem & 1680 & 54 & 3.21\% \\
B & it may be asked & 1680 & 0 & 0.00\% \\
B & one may suppose & 1680 & 15 & 0.89\% \\
B & any hedge phrase & 1680 & 340 & 20.24\% \\
\hline
\end{tabular}
\end{table*}

\newpage
\newpage
% \bibliographystyle{plain}
\bibliographystyle{unsrt}
\bibliography{references}

\end{document}